\documentclass[10pt]{beamer}
\usepackage{polski}
\usepackage[utf8]{inputenc}
\usepackage[T1]{fontenc}
\usepackage{carlito}
\usepackage{tikz}
\usetikzlibrary{positioning}
\usepackage{graphicx}
\usepackage{booktabs}

% Wszystkie zdjęcia (także tła szablonu NewPwr) leżą w photos/
\graphicspath{{photos/}}

%=============== Ustawienia szablonu ===============
\usetheme[horizontal=true, hr=false]{NewPwr}
\setbeamertemplate{footline}{%
  \leavevmode%
  \hbox{%
    \hspace{10pt}%
    \begin{beamercolorbox}[wd=\paperwidth,ht=2.5ex,dp=1.125ex]{page number in head/foot}%
      \usebeamerfont{page number in head/foot}\insertframenumber%
    \end{beamercolorbox}%
  }%
  \vskip10pt%
}
\setbeamercolor{page number in head/foot}{fg=white}
\setbeamerfont{page number in head/foot}{size=\Large}
\setbeamertemplate{navigation symbols}{}

%=============== Dane Tytułowe ===============
\title[]{Tytuł prezentacji}
\author[Koło Naukowe]{\textbf{} \\
    Autor 1 \\
    Autor 2}
\date{\today}

\begin{document}

% --- Strona tytułowa ---
\begin{frame}
\titlepage
\end{frame}

% --- Plan prezentacji ---
\begin{frame}{Plan prezentacji}
    \begin{enumerate}
        \item \textbf{Część 1} – opis
        \item \textbf{Część 2} – opis
        \item \textbf{Część 3} – opis
    \end{enumerate}
\end{frame}

\section{Część 1}

% --- Blok + zdjęcie ---
\begin{frame}{Slajd z blokiem i zdjęciem}
    \begin{block}{Nagłówek bloku}
        Treść bloku.
    \end{block}
    \begin{figure}
        \centering
        \includegraphics[width=0.6\linewidth]{example-image}
        \caption{Podpis zdjęcia}
    \end{figure}
\end{frame}

% --- Dwie kolumny ---
\begin{frame}{Slajd z dwiema kolumnami}
    \begin{columns}[T]
        \begin{column}{0.48\textwidth}
            \textbf{Lewa kolumna}
            \begin{itemize}
                \item Punkt 1
                \item Punkt 2
            \end{itemize}
        \end{column}
        \begin{column}{0.48\textwidth}
            \textbf{Prawa kolumna}
            \includegraphics[width=\linewidth]{example-image}
        \end{column}
    \end{columns}
\end{frame}

% --- Tabela ---
\begin{frame}{Slajd z tabelą}
    \begin{table}
        \centering
        \begin{tabular}{lll}
            \toprule
            Kolumna 1 & Kolumna 2 & Kolumna 3 \\
            \midrule
            a & b & c \\
            d & e & f \\
            \bottomrule
        \end{tabular}
        \caption{Podpis tabeli}
    \end{table}
\end{frame}

\section{Podsumowanie}

\begin{frame}{Podsumowanie}
    \begin{block}{Omówione zagadnienia}
        \begin{itemize}
            \item Punkt 1
            \item Punkt 2
        \end{itemize}
    \end{block}
\end{frame}

\begin{frame}{Dziękuję za uwagę}
    \begin{center}
        \Huge \textbf{Pytania?}
    \end{center}
\end{frame}

\end{document}
